\chapter{Edge AI System Implementation}
\label{chap:deployment}

In this chapter, the complete embedded implementation of the Edge AI grape leaf disease detection system on ESP32-S3 microcontroller is presented. Following the chosen models design and development described in Chapter~\ref{chap:system_architecture}, this chapter addresses the critical steps of deploying two neural network models on a resource-constrained microcontroller particularly ESP32-S3.


This chapter begins with a description of the chosen hardware and its development framework, including the ESP-IDF development framework,
and the ESP-DL inference engine presented in Section~\ref{sec:hw_platform}. Section~\ref{sec:quantization_pipeline} details the procedures of post-training quantization pipeline that converts both models from floating point to INT8 integer representations compatible with the ESP-DL inference engine. Section~\ref{sec:firmware_memory}  discusses firmware architecture and memory management strategies, including challenges related to model storage, initialization sequence, and buffer allocation. Section~\ref{sec:inference_pipeline} demonstrates complete inference pipeline implementation, from image capturing through inference results aggregation. Finally, Section~\ref{sec:implementation_constraints} documents the key implementation constraints and design trade-offs.

\section{Hardware and Development Framework}
\label{sec:hw_platform}

\subsection{ESP32-S3 Microcontroller Architecture}
\label{subsec:esp32s3_architecture}
The deployment target is ESP32-S3 system-on-chip (SoC) module manufactured by Espressif Systems~\cite{esp32s3_datasheet}. The ESP32-S3 integrates a dual-core Xtensa LX7 processor operating at 240~MHz, providing computational capabilities suitable for embedded deep learning inference while maintaining low power consumption. The Xtensa LX7 cores support single-instruction multiple-data (SIMD) operations through the PIE (Processor Interface Extension) instruction set, enabling efficient execution of vectorized operations common in neural network inference. The ESP-DL inference engine leverages these SIMD capabilities. The SoC includes a 512~KB internal SRAM, 8~MB external PSRAM, and 16~MB external flash memory. The memory architecture is critical for model storage and inference execution:

\begin{itemize}
    \item \textbf{Internal SRAM (512~KB):} High-speed on-chip memory with single-cycle access latency, partitioned between instruction RAM (IRAM), data RAM (DRAM), and DMA-accessible regions. The ESP-DL runtime uses SRAM for frequently accessed weight tiles via a least-recently-used (LRU) caching policy, storing up to 64~KB of active parameters.
    \item \textbf{External PSRAM (8~MB):} Octal SPI PSRAM operating at 80~MHz with Hybrid Wrap burst mode and 32-byte burst length. PSRAM provides the primary storage for model, weights and activations buffer, camera frame buffers, and inference intermediate tensors. The ESP-DL runtime manages PSRAM allocations through the ESP-IDF heap allocator, with explicit control over memory regions via capability flags.
    \item \textbf{External Flash (16~MB):} QIO-mode (Quad Input/Output) flash storage at 80~MHz clock speed, used for firmware binary storage including model weights embedded in the \texttt{.rodata} section. The flash partition table allocates a 4~MB factory application partition. The ESP-DL runtime accesses model weights directly from flash via memory-mapped I/O, caching active tiles in SRAM to mitigate flash access latency.
\end{itemize}
In this deployment, I used Freenove ESP32-S3-WROOM~1 development board~\cite{freenove2024esp32s3} with OV3660 camera sensor connected via 8-bit parallel DVP interface. The OV3660 captures at VGA resolution (640$\times$480) with on-chip JPEG compression, supporting clock frequencies up to 40~MHz~\cite{esp32s3_datasheet}. Camera data is transferred to PSRAM through DMA channels. The ESP32-S3's integrated Wi-Fi and Bluetooth capabilities are not utilized in this deployment.
\subsection{ESP-IDF Development Framework}
\label{subsec:esp_idf}

The firmware developed using Espressif IoT Development Framework (ESP-IDF version 5.3.3), Espressif's official C/C++ development platform for ESP32-series microcontrollers. ESP-IDF provides a FreeRTOS-based real-time operating system with hardware abstraction layers for peripheral management, memory allocation, and power management. The build system uses CMake with Ninja, compiling the application into a single monolithic binary that includes the bootloader, partition table, and application code. The ESP-IDF framework offers the following features relevant to this deployment:
\begin{itemize}
    \item \textbf{Heap allocator:} Manages both internal SRAM and external PSRAM through a unified \texttt{heap\_caps} API. Memory allocations can be directed to specific memory regions via capability flags (e.g., \texttt{MALLOC\_CAP\_SPIRAM} for PSRAM allocation, \texttt{MALLOC\_CAP\_DMA} for DMA-accessible SRAM).
    \item \textbf{Camera driver:} Configures the DVP camera interface, manages DMA transfers, and provides double-buffered frame capture. The driver supports multiple pixel formats including JPEG, RGB565, and YUV422.
    \item \textbf{SPI flash driver:} provides low-level access to external flash memory, including support for memory-mapped I/O. The \texttt{esp\_partition} API allows defining custom flash partitions for model storage, while the \texttt{esp\_spi\_flash} API enables direct reads from flash with optional caching.
    \item \textbf{Power management:} ESP-IDF's power management API allows configuring sleep modes and dynamic frequency scaling to optimize energy consumption during idle periods. This enables the duty-cycled operation of inference pipeline, to extend battery life.
\end{itemize}

\subsection{ESP-DL Inference Engine}
\label{subsec:esp_dl}

ESP-DL~\cite{esp_dl} is Espressif's neural network inference library optimized for ESP32-series microcontrollers. The library provides a C++ inference runtime that executes quantized INT8 models directly on the Xtensa LX7 cores, leveraging SIMD instructions for vectorized computation. ESP-DL supports a range of layer types commonly used in convolutional neural networks, including convolution, depthwise convolution, fully connected, activation functions, pooling, and softmax. The library includes optimized kernels for the ESP32-S3's architecture, with specific attention to memory access patterns and cache utilization. The ESP-DL runtime manages model loading, memory allocation for weights and activations, and execution scheduling. Key features relevant to this deployment include:

\begin{itemize}
    \item \textbf{Model loading:} (\texttt{dl::Model}) class that loads quantized models from flash RODATA, with an optional \texttt{minimize()} method that deletes unused variables to reduce memory footprint at the cost of debug capability.
    \item \textbf{Image preprocessing:}  (\texttt{dl::image::ImagePreprocessor}) class that handles input image transformations required by the models, including resizing, normalization, and color space conversion. The preprocessor supports various pixel formats; for this deployment, input images are processed in RGB888 format for hardware compatibility model input requirements.
    \item \textbf{Detection post-processing:} (\texttt{dl::detect::ESPDetPostProcessor}) class that uses the raw output tensors from the detection model to produce final bounding box predictions. The post-processor implements non-maximum suppression (NMS), confidence thresholding, and coordinate transformation for object detection models, handling INT8 dequantization internally.
    \item \textbf{Memory management:} The runtime allocates activation buffers in PSRAM and manages an LRU cache of weight tiles in SRAM to optimize inference performance while respecting memory constraints. The library provides APIs for explicit control over memory allocation and deallocation, enabling the zero-copy buffer reuse strategy described in Section~\ref{subsec:buffer_allocation}.
    \item \textbf{Software JPEG decoder:} (\texttt{dl::image::sw\_decode\_jpeg()}) function that decodes JPEG-compressed camera frames to RGB888 format, as the ESP32-S3 lacks a dedicated hardware JPEG decoder.

\end{itemize}


\section{Model Quantization and Format Conversion}
\label{sec:quantization_pipeline}

Deploying neural network models on microcontroller devices requires converting floating-point to INT8 representations compatible with target inference engine. This section describes the quantization pipeline that conversion of both models ESPDet-Pico and MobileNetV2 from PyTorch FP32 to the ESPDL INT8 format. The conversion pipeline follows the model quantization techniques in Espressif's official ESP-DL deployment documentation ~\cite{esp_dl_mobilenetv2_tutorial}, adapted to the specific requirements of each model architecture. 

\subsection{The ESP-PPQ Quantization Framework}
\label{subsec:esp_ppq}

Both models quantized using ESP-PPQ~\cite{esp_ppq2023}, Espressif's Post-Processing Quantization framework. ESP-PPQ is extended over the open-source PPQ library with hardware-specific quantization configurations targeting ESP-DL inference runtime. The framework applies per-tensor symmetric quantization with ($s = 2^e$) factors, where floating point values mapped to the 8-bit integer range $[-128, 127]$:

\begin{equation}
x_{\text{int8}} = \text{clamp}\left(\text{round}\left(\frac{x_{\text{float}}}{s}\right), -128, 127\right)
\label{eq:quantization}
\end{equation}

where $x_{\text{float}}$ is the original floating point value, $s$ is the scale factor, and $\text{clamp}$ to ensure the quantized value fits within the 8-bit signed integer range. The quantization exponent $e$ computed during calibration to minimize the quantization error over its representative dataset. The power-of-2 constraint on the scale factor enables efficient dequantization via bit-shift operation on the Xtensa LX7 processor, avoiding costly floating-point operation during inference.

The ESP specific quantization setting \texttt{QuantizationSettingFactory.espdl\_setting} with the \texttt{ESPDL\_INT8} target applies a layerwise equalization to distribute weights and activations across consecutive layers, reducing the dynamic range that must be with in 8-bit precision. This strategy particularly important for convolutional layers with highly skewed weight distributions. Thus, this addresses the challenge of MobileNetV2's depthwise separable convolutions, where weight distribution varies significantly across channels.

The ESP-PPQ framework performs quantization first, it collects activation statistics from the calibration dataset to determine optimal scale factors $s$ for each layer. Then, it applies the quantization conversion weights and activations, converting the model to INT8 format. Finally, the quantized model exported in a \texttt{.espdl} binary format compatible with the ESP-DL runtime. 




\subsection{The ESPDet-Pico Model Quantization Pipeline}
\label{subsec:espdet_quantization}
After training the ESPDet-Pico leaf localization model in section~\ref{sec:espdet_pico}, its PyTorch weights were exported into ONNX (Open Neural Network Exchange) format to obtain a standardized intermediate representation. The export used a fixed input shape of [1, 3, 320, 416] (NCHW) and opset 13 to ensure compatibility with the ESP-DL pipeline. The ONNX model was simplified using \texttt{onnxsim} to eliminate redundant operations and fuse batch normalization layers, resulting in a more efficient graph for quantization. Then the ESP-PPQ quantization was applied to the simplified ONNX model with 32 batches of calibration data and equalization enabled. This achieved 73\% reduction in model size. Finally, the quantized model was converted to \texttt{.espdl} format for deployment on the ESP32-S3, enabling real-time leaf localization. In Figure~\ref{fig:espdet_quantization_pipeline} illustrated the complete quantization pipeline.

\begin{figure}[htbp]
\centering
\begin{tikzpicture}[
    node distance=0.5cm,
    box/.style={
        rectangle, rounded corners=2mm,
        draw=NavyBlue, fill=lightblue,
        text width=13.5cm, align=left,
        minimum height=1.25cm,
        drop shadow
    },
    arrow/.style={-{Stealth[length=3mm]}, thick, NavyBlue}
]

\node[box] (data) {\textbf{1. Dataset Preparation (Roboflow)}\\
\hspace{0.5cm}513 images (328 train + 82  val + 103 test) | YOLOv11 format \\ \hspace{0.5cm} | Single class: `leaf'};

\node[box, below=of data] (train) {\textbf{2. Model Training (ESPDet-Pico)}\\
\hspace{0.5cm}677 epochs | AdamW optimizer | Precision: 64.8\%, mAP50: 56.4\%};

\node[box, below=of train] (onnx) {\textbf{3. Convert model inot ONNX Format}\\
\hspace{0.5cm}PyTorch $\rightarrow$ ONNX (FP32, 1.06 MB) | ESP-specific modifications};

\node[box, below=of onnx] (quant) {\textbf{4. INT8 Quantization (ESP-PPQ)}\\
\hspace{0.5cm}32 calibration batches | Equalization enabled};

\node[box, below=of quant] (deploy) {\textbf{5. Prepare ESP-DL Compatible Model (.espdl)}\\
\hspace{0.5cm}478 KB | INT8 | ESP32-S3 ready | Inference: $\approx$360 ms/frame};

\draw[arrow] (data) -- (train);
\draw[arrow] (train) -- (onnx);
\draw[arrow] (onnx) -- (quant);
\draw[arrow] (quant) -- (deploy);

\end{tikzpicture}
\caption{ESPDet-Pico quantization pipeline from dataset preparation to ESP32-S3 deployment-ready INT8 model.}
\label{fig:espdet_quantization_pipeline}
\end{figure}

\subsubsection*{The Simplified ONNX Format}

Pretrained PyTorch ESPDet-Pico model converted to ONNX format using a custom export path that produces six separate output tensors rather than single concatenated tensor used in standard YOLO inference. A custom \texttt{ESP\_YOLO} wrapper class replaces the standard detection head's forward method with \texttt{export\_onnx\_forward}, which outputs separate box regression and classification score tensors for each of the three detection scales:

\begin{equation}
\text{Output} = \{(\mathbf{B}_s, \mathbf{S}_s)\}_{s \in \{P3, P4, P5\}}
\end{equation}

where $\mathbf{B}_s \in \mathbb{R}^{1 \times 4 \times H_s \times W_s}$ contains raw box edge distances and $\mathbf{S}_s \in \mathbb{R}^{1 \times 1 \times H_s \times W_s}$ contains class logits at scale $s$. This six-tensor format avoids complex tensor concatenation and reshaping operations that are computationally expensive on the ESP32-S3 and allows the ESP-DL post-processor to handle each scale independently.


\subsubsection*{INT8 Quantization}

The simplified ONNX model is quantized following the procedures of ESP-DL model quantization guide~\cite{esp_dl_mobilenetv2_tutorial}. It needs a calibration dataset to collect activation statistics to determine quantization parameters. The crital constraint for the ESPDet-Pico model was its rectangular input shape ($320 \times 416$), which requires a custom calibration pipeline. And the the batch size must be set to 1 during calibration and proceeds for 32 steps with equalization enabled. 

The qunatization and format conversion can be performed by calling a single function from the ESP-PPQ module, which internally performs INT8 calibration and exports the result directly to the \texttt{.espdl} binary format. Then detection model's INT8 weights become total 478~KB, a 55\% reduction from the 1.06~MB FP32 model.

\subsection{The MobileNetV2 Model Quantization Pipeline}
\label{subsec:mobilenet_quantization}
To deploy the MobileNetV2 disease classification model on the ESP32-S3, the trained network converted into ONNX format and then quantized to \textbf{INT8} using \textbf{ESP-PPQ}. I used similar quantization configuration as the ESPDet-Pico model, ensuring ESP-DL compatibility. The same steps as ESPDet-Pico, with adjustments for the model's input shape ($128 \times 128$) and architecture. The complete quantization pipeline is illustrated in Figure~\ref{fig:mobilenet_quantization_pipeline}.

\begin{figure}[htbp]
\centering
\begin{tikzpicture}[
    node distance=0.5cm,
    box/.style={
        rectangle, rounded corners=2mm,
        draw=NavyBlue, fill=lightblue,
        text width=13.5cm, align=left,
        minimum height=1.25cm,
        drop shadow
    },
    arrow/.style={-{Stealth[length=3mm]}, thick, NavyBlue}
]

\node[box] (data) {\textbf{1. Dataset Preparation }\\
\hspace{0.5cm}4062 images (2,761 train + 487  val + 814 test) | Labeled Images \\ \hspace{0.5cm} | four classes: `Black Rot', `Esca', `Leaf Blight', `Healthy'};

\node[box, below=of data] (train) {\textbf{2. MobileNetV2 Fine-tune }\\
\hspace{0.5cm} 50 epochs | AdamW optimizer | Precision: 99.8\%};

\node[box, below=of train] (onnx) {\textbf{3. Convert model into ONNX Format}\\
\hspace{0.5cm}PyTorch $\rightarrow$ ONNX (FP32, 8.48 MB) | ESP-specific modifications};

\node[box, below=of onnx] (quant) {\textbf{4. INT8 Quantization (ESP-PPQ)}\\
\hspace{0.5cm}32 calibration batches | Equalization enabled};

\node[box, below=of quant] (deploy) {\textbf{5. Prepare ESP-DL Compatible Model (.espdl)}\\
\hspace{0.5cm}2.27 MB | INT8 | ESP32-S3 ready | Precision: 99.5\%};

\draw[arrow] (data) -- (train);
\draw[arrow] (train) -- (onnx);
\draw[arrow] (onnx) -- (quant);
\draw[arrow] (quant) -- (deploy);

\end{tikzpicture}
\caption{MobileNetV2 quantization pipeline from dataset preparation to ESP32-S3 deployment-ready INT8 model}
\label{fig:mobilenet_quantization_pipeline}
\end{figure}



\subsection{Quantization Accuracy Preservation}
\label{subsec:quantization_accuracy}

Post-training quantization introduces rounding errors that can degrade model accuracy. The layerwise equalization strategy is selected specifically to address this challenge. To contextualize the quantization accuracy trade-off, Espressif's quantization documentation~\cite{esp_dl_mobilenetv2_tutorial} reports comparative results for three strategies applied to a standard MobileNetV2 model (ImageNet classification, 1000 classes):

\begin{itemize}
    \item \textbf{Default 8-bit quantization:} Produces 25.6\% last-layer error, reducing Top-1 accuracy from 71.9\% (float) to 60.5\%.
    \item \textbf{Mixed precision (INT16 for worst layer):} Reduces last-layer error to 9.1\%, achieving 69.6\% accuracy.
    \item \textbf{Layerwise equalization:} Achieves 9.0\% last-layer error with 69.8\% accuracy, matching mixed precision while maintaining uniform INT8 quantization.
\end{itemize}

The equalization approach is adopted for both models because it preserves uniform INT8 quantization while achieving accuracy comparable to more complex strategies. The grape disease MobileNetV2 model maintains test accuracy above 99.5\% after INT8 quantization, with less than 0.3 percentage point degradation from the FP32 baseline. This robustness reflects MobileNetV2's batch normalization layers and the inherently quantization-friendly nature of inverted residual blocks~\cite{jacob2018quantization}.


\section{Firmware Architecture and Memory Management}
\label{sec:firmware_memory}

Deploying two neural network models on a microcontroller with 8~MB PSRAM, 16~MB flash RAM, and 512~KB SRAM imposes careful memory management requirements. This section presents the firmware architecture decisions governing model storage, initialization sequencing, and buffer allocation strategies, each of which critically influences both deployment reliability and inference performance.
\subsection{ Model Storage and Loading Strategy}
A key design decision in the firmware architecture is the strategy for storing and loading model weights into the inference engine. Following training, INT8 quantization, conversion to \texttt{.espdl} format, the neural network binaries are stored in the flash memory of the ESP32-S3, which is non-volatile and has a capacity of 16~MB. The firmware is designed to load these models into PSRAM at runtime for inference execution as shown in Figure~\ref{fig:model_storage_strategy}.


\begin{figure}[htbp]
\centering

\begin{tikzpicture}[
    scale=0.85,
    transform shape,
    box/.style={rectangle, draw, rounded corners=2mm, minimum width=3.2cm, minimum height=0.8cm, align=center, font=\small},
    smallbox/.style={rectangle, draw, rounded corners=2mm, minimum width=3.2cm, minimum height=0.7cm, align=center, font=\small},
    container/.style={rectangle, draw, rounded corners=3mm, thick, inner sep=8pt},
    arrow/.style={->, >=stealth, thick}
]

% =========================
% Flash container + models
% =========================
\node[box, fill=orange!20] (flash) {\textbf{Flash Memory}\\16 MB};

\node[smallbox, fill=blue!20, below=0.6cm of flash] (det) {\textbf{ESPDet-Pico}\\478 KB};
\node[smallbox, fill=blue!20, below=0.35cm of det] (cls) {\textbf{MobileNetV2 (INT8)}\\2.27 MB};

\node[container, fit=(flash)(det)(cls)] (flashbox) {};

% Flash label (aligned cleanly)
\node[font=\scriptsize\bfseries, fill=white, inner sep=2pt, anchor=west]
    at ($(flashbox.north west)+(0.1,0.25)$) {Non-volatile Storage};

% =========================
% PSRAM container + runtime buffers
% =========================
\node[box, fill=green!20, right=4.6cm of det] (psram) {\textbf{PSRAM}\\8 MB};

\node[smallbox, fill=yellow!20, below=0.45cm of psram] (buffers)
{\textbf{Runtime Buffers}\\Input / Output / Activations};

\node[container, fit=(psram)(buffers)] (psrambox) {};

% PSRAM label (aligned cleanly)
\node[font=\scriptsize\bfseries, fill=white, inner sep=2pt, anchor=west]
    at ($(psrambox.north west)+(0.1,0.25)$) {Volatile Runtime Memory};

% =========================
% Arrows (data movement)
% =========================
\draw[arrow] (det.east) -- node[above, font=\scriptsize] {load weights} (psram.west);
\draw[arrow] (cls.east) -- node[below, font=\scriptsize] {load weights} (psram.west);
\draw[arrow] (psram) -- (buffers);

\end{tikzpicture}

\vspace{0.25cm}

\begin{itemize}\itemsep2pt
    \item Both models packed into a single binary (\texttt{grape\_leaf\_detect.espdl}, 2.73~MB total) stored in flash memory
    \item PSRAM reserved for runtime tensors, activation buffers, and camera framebuffers
    \item Models accessed directly from flash via memory-mapped I/O with SRAM caching
\end{itemize}

\caption{\small Model storage and loading strategy for ESP32-S3 deployment. Both ESPDet-Pico and MobileNetV2 models are embedded in flash memory and loaded into PSRAM at runtime for inference execution.}
\label{fig:model_storage_strategy}
\end{figure}

\subsection*{Model Storage Strategies}
\label{subsec:model_storage}

The strategy for loading model weights into the inference engine significantly affects initialization latency, runtime reliability, and PSRAM availability. Three distinct approaches were evaluated during firmware development:

\textbf{1. Dedicated Flash Partition Loading:} Models stored in custom flash partitions at a fixed offset and loaded via \texttt{MODEL\_LOCATION\_IN\_FLASH\_PARTITION}. This approach introduces cache coherency issues and memory management unit (MMU) mapping conflicts when accessing partition data, resulting in non-deterministic model loading failures that preclude reliable autonomous deployment.

\textbf{2. Memory-Mapped Access:} Direct memory-mapped I/O with, reading model parameters directly from flash without copying into PSRAM. Although this approach minimizes memory footprint, it exhibits instability under high memory utilization and cache lines holding model parameters are thrown out when the camera DMA controller claims buffer regions, causing inference failures.

\textbf{3. RODATA Embedding with Direct Flash Access (adopted):} Both models compiled into the firmware binary as constant arrays in the \texttt{.rodata} section. During the build process, the ESP-IDF build system uses the CMake function to convert the \texttt{.espdl} model files into object files linked directly into the final firmware image. The GCC linker automatically generates special symbols marking the memory addresses where each embedded binary begins and ends, following the naming convention \texttt{\_binary\_<filename>\_start} and \texttt{\_binary\_<filename>\_end}.  Models remain in flash and are accessed directly through the instruction cache without copying to PSRAM. The ESP-DL runtime internally caches frequently accessed weight tiles in SRAM via an LRU eviction policy (up to 64~KB), balancing flash access latency with memory capacity.



The RODATA approach reduces initialization time for both models combined, while conserving PSRAM exclusively for activation buffers and camera framebuffers. Hence, this strategy was adopted for the production firmware, ensuring stable model loading and maximizing available PSRAM for runtime operations.



\subsection{Initialization Order of Camera and Models}
\label{subsec:init_order}
One of the most challenging issues faced during firmware deployment was the initialization order between the camera module and the neural network models. The ESP32-S3's PSRAM is shared between the camera frame buffers and the model activation buffers, and its allocation strategy can lead to fragmentation that prevents successful model loading if not managed correctly. If the camera is initialized first, it allocates large contiguous buffers for frame capture, which can fragment the PSRAM heap and block subsequent allocations for model weights and activations. Conversely, if the models are initialized first, they can secure their required contiguous memory regions before the camera's DMA buffers are allocated, allowing the system to function correctly despite fragmentation caused by the camera. This critical dependency was identified through systematic testing of different initialization sequences and monitoring PSRAM usage at each step. The final adopted sequence is as follows, illustrated in  Figure~\ref{fig:initialization_sequence}:

\begin{figure}[htbp]
\centering
\begin{tikzpicture}[
    node distance=0.8cm,
    stage/.style={
        rectangle, rounded corners=2mm,
        draw=NavyBlue, fill=lightblue,
        text width=5.5cm, align=left,
        minimum height=1.8cm,
        drop shadow,
        font=\small
    },
    membar/.style={
        rectangle,
        draw=black,
        minimum height=0.35cm,
        anchor=west
    },
    arrow/.style={-{Stealth[length=2mm]}, thick, NavyBlue}
]

% Timeline stages
\node[stage] (stage1) at (0,0) {
    \textbf{Step 1: Leaf Detection Model}\\
    \textbf{Time:} 144~ms\\
    \textbf{PSRAM used:} 1,095~KB\\
    \textbf{PSRAM free:} 7,095~KB
};

\node[stage, below=of stage1] (stage2) {
    \textbf{Step 2: Classification Model}\\
    \textbf{Time:} 348~ms\\
    \textbf{PSRAM used:} +2,888~KB\\
    \textbf{PSRAM free:} 4,207~KB
};

\node[stage, below=of stage2] (stage3) {
    \textbf{Step 3: Camera Driver}\\
    \textbf{Time:} --\\
    \textbf{PSRAM used:} $\approx$180~KB\\
    \textbf{Runtime free:} $\approx$3,910~KB
};

\draw[arrow] (stage1.south) -- (stage2.north);
\draw[arrow] (stage2.south) -- (stage3.north);

\def\memscale{0.0008}  % scale factor: 1 KB = 0.0008 cm

\node[font=\scriptsize\bfseries, right=1.5cm of stage1.north east, anchor=south west] (mem0label) 
    {Initial: 8,190~KB};

% Stage 1 memory bar
\node[membar, fill=red!40, minimum width=0.5cm, right=1.5cm of stage1.east] (mem1used) {};
\node[membar, fill=green!30, minimum width=3.2cm, right=0cm of mem1used.east] (mem1free) {};
\node[font=\scriptsize, right=0.05cm of mem1free.east] {7,095~KB};

% Stage 2 memory bar
\node[membar, fill=red!40, minimum width=1.8cm, right=1.5cm of stage2.east] (mem2used) {};
\node[membar, fill=green!30, minimum width=1.9cm, right=0cm of mem2used.east] (mem2free) {};
\node[font=\scriptsize, right=0.05cm of mem2free.east] {4,207~KB};

% Stage 3 memory bar (fragmented)
\node[membar, fill=red!40, minimum width=1.9cm, right=1.5cm of stage3.east] (mem3used) {};
\node[membar, fill=green!40, minimum width=1.8cm, right=0cm of mem3used.east] (mem3free) {};
\node[font=\scriptsize, right=0.05cm of mem3free.east, text width=1.8cm, align=left] {\textasciitilde3.9~MB};

% Legend
\node[font=\scriptsize, below=0.8cm of stage3.south west, anchor=north west] (legend) {
    \textbf{Legend:} 
    \colorbox{red!40}{\phantom{x}} Used \quad
    \colorbox{green!30}{\phantom{x}} Free \quad
};

\end{tikzpicture}
\caption{Initialization sequence and PSRAM utilization for model loading and camera initialization.}
\label{fig:initialization_sequence}
\end{figure}


\subsection{Buffer Allocation and Zero-Copy Reuse Strategy}
\label{subsec:buffer_allocation}

To minimize runtime memory allocation overhead and prevent heap fragmentation during continuous operation, the inference pipeline employs a \textbf{single-allocation, zero-copy buffer strategy}:

\begin{itemize}
    \item \textbf{Classification input buffer (49~KB):} A single $128 \times 128 \times 3$ RGB888 buffer allocated in PSRAM during classifier initialization and reused across all detected leaf regions within a frame. This eliminates repeated \texttt{malloc}/\texttt{free} calls that would accumulate heap fragmentation over hours of operation.
    \item \textbf{Camera framebuffers (122~KB):} Two framebuffers allocated by the camera driver in PSRAM via \texttt{CAMERA\_FB\_IN\_PSRAM}. The driver manages double-buffering internally: one buffer captures the next frame while the other undergoes processing.
    \item \textbf{JPEG decode buffer ($\sim$900~KB):} The decoded RGB888 frame ($640 \times 480 \times 3 = 921{,}600$ bytes) is allocated dynamically by the software JPEG decoder. This temporary allocation is freed immediately after detection inference completes and leaf regions are extracted from the decoded buffer.
    \item \textbf{Model activation buffers ($\sim$2~MB total):} Intermediate feature maps allocated by the ESP-DL memory manager during model initialization. These buffers persist throughout execution using a static allocation plan computed offline during quantization.
\end{itemize}

The elimination of dynamic allocation in the inference loop ensures deterministic memory behavior and prevents the gradual heap degradation that would compromise long-term deployment stability.


\section{Inference Pipeline Implementation}
\label{sec:inference_pipeline}

The complete inference pipeline orchestrates five sequential stages: camera acquisition, JPEG decoding, leaf detection, region extraction with preprocessing, and disease classification, followed by cross-leaf aggregation for frame-level diagnosis.

\subsection{Camera Acquisition and JPEG Decoding}
\label{subsec:camera_acquisition}

Image acquisition begins with the OV3660 camera module configured for VGA resolution ($640 \times 480$) with on-chip JPEG compression. Camera configuration prioritizes deployment robustness:

\begin{itemize}
    \item \textbf{Pixel format:} \texttt{PIXFORMAT\_JPEG} with quality setting 12 (scale 0--63, lower = higher quality), balancing compression ratio ($\sim$10:1) with minimal blocking artifacts that could degrade detection accuracy. Captured JPEG frames occupy typically 12--80~KB.
    \item \textbf{Automatic exposure and white balance:} Enabled via sensor control registers (\texttt{set\_exposure\_ctrl}, \texttt{set\_whitebal}, \texttt{set\_awb\_gain}) to adapt to varying vineyard lighting conditions without manual recalibration.
    \item \textbf{Frame buffering:} Dual-buffer mode (\texttt{fb\_count=2}) enables capture-while-processing: the camera populates buffer 1 while the inference pipeline processes buffer 0.
\end{itemize}

Captured JPEG frames are decoded to RGB888 format using ESP-DL's software JPEG decoder (\texttt{dl::image::sw\_decode\_jpeg()}), as the ESP32-S3 lacks hardware JPEG decoding. The decoder writes directly to a heap-allocated RGB888 buffer ($640 \times 480 \times 3 = 921{,}600$ bytes) in PSRAM, avoiding intermediate color space conversions.

\subsection{Leaf Detection and Coordinate Mapping}
\label{subsec:leaf_detection_deployment}

The decoded RGB888 frame undergoes ESPDet-Pico detection inference to localize individual leaf regions. The \texttt{dl::image::ImagePreprocessor} resizes the $640 \times 480$ input to $416 \times 320$ via letterbox padding with gray fill (RGB$=[114, 114, 114]$), preserving the original aspect ratio. Input pixels are normalized by scaling to $[0, 1]$ (mean$=[0,0,0]$, std$=[255,255,255]$), consistent with the model's training preprocessing.

Detection inference produces six output tensors (box and score predictions at three scales), processed by the \texttt{ESPDetPostProcessor} with the following parameters:

\begin{itemize}
    \item \textbf{Confidence threshold:} $\tau = 0.25$, tuned empirically to balance recall and computational cost.
    \item \textbf{NMS IoU threshold:} 0.7, eliminating duplicate detections of overlapping leaves.
    \item \textbf{Maximum detections:} 10 per frame, limiting worst-case classification latency.
    \item \textbf{Anchor strides:} $\{8, 16, 32\}$ for the three detection scales, with grid sizes corresponding to P3/8, P4/16, and P5/32 feature maps.
\end{itemize}

Detected bounding box coordinates are mapped from the $416 \times 320$ detection space back to the original $640 \times 480$ frame space via inverse letterbox transformation:
\begin{equation}
x_{\text{orig}} = \frac{(x_{\text{det}} - \text{pad}_x) \times W_{\text{orig}}}{W_{\text{det}} - 2\text{pad}_x}, \quad
y_{\text{orig}} = \frac{(y_{\text{det}} - \text{pad}_y) \times H_{\text{orig}}}{H_{\text{det}} - 2\text{pad}_y}
\end{equation}
This transformation ensures that cropped leaf regions align precisely with the original RGB888 frame for downstream classification.

\subsection{Region Cropping and Classification Preprocessing}
\label{subsec:region_cropping}

For each detected bounding box, the corresponding leaf region is extracted from the $640 \times 480$ RGB888 frame and prepared for MobileNetV2 classification. The cropping operation implements \textbf{square expansion with letterbox padding}: the rectangular bounding box is expanded to a square centered on the box's midpoint, with side length equal to $\max(\text{width}, \text{height})$. This preserves the leaf's aspect ratio and prevents distortion that could impair classification accuracy on elongated or irregularly shaped leaves.

The squared region is resized to $128 \times 128$ pixels using \textbf{nearest-neighbor interpolation}:
\begin{equation}
I_{128}[y, x] = I_{\text{frame}}\!\left[\left\lfloor y \cdot s_y \right\rfloor,\; \left\lfloor x \cdot s_x \right\rfloor\right]
\end{equation}
where $s_x = W_{\text{bbox}} / 128$ and $s_y = H_{\text{bbox}} / 128$. Nearest-neighbor interpolation is chosen over bilinear interpolation to minimize per-crop preprocessing time (1--2~ms versus 10--12~ms per crop), a significant saving when processing up to 10 leaves per frame.

Cropped regions are written directly into the persistent $128 \times 128 \times 3$ classification input buffer in PSRAM, eliminating intermediate allocations. Pixels outside the original frame bounds (due to square expansion) are filled with gray padding (RGB$=[127, 127, 127]$), approximating ImageNet's mean pixel value.

\subsection{Disease Classification and INT8 Dequantization}
\label{subsec:classification_dequant}

The cropped $128 \times 128$ RGB888 leaf image undergoes MobileNetV2 classification inference. The model's input tensor pointer is set directly to the pre-allocated classification buffer (\texttt{input->set\_element\_ptr(input\_buffer)}), and inference is triggered via \texttt{model->run()}, executing the quantized INT8 forward pass through 17 inverted residual blocks.

The model outputs a 4-dimensional INT8 logit vector representing unnormalized class scores for \texttt{Black\_rot} (index 0), \texttt{Esca} (index 1), \texttt{Healthy} (index 2), and \texttt{Leaf\_blight} (index 3), ordered alphabetically per PyTorch's \texttt{ImageFolder} class indexing convention. The INT8 logits must be dequantized to floating-point before probability interpretation:

\begin{equation}
\text{logit}_{\text{float}}^{(c)} = \text{logit}_{\text{int8}}^{(c)} \times 2^{e}
\label{eq:dequant}
\end{equation}

where $e$ denotes the quantization exponent embedded in the output tensor's metadata, retrieved via \texttt{output->exponent}. The scale factor $s = 2^e$ is computed using the ESP-DL macro \texttt{DL\_SCALE(e)}.

A significant implementation constraint arises from the asymmetry in how the ESP-DL framework handles dequantization across model types. The detection model's \texttt{ESPDetPostProcessor} performs dequantization internally as part of its NMS pipeline, transparently converting INT8 outputs to floating-point coordinates and confidence scores. In contrast, the classification model outputs \textbf{raw INT8 logits} that require explicit dequantization before softmax computation. Omitting this dequantization step --- feeding raw INT8 values (range $[-128, 127]$) directly to the softmax function --- causes extreme probability concentration due to the large magnitude differences, resulting in near-100\% confidence on a single class regardless of input (Section~\ref{subsec:class_mapping}).

Dequantized logits are converted to probabilities via numerically stable softmax:
\begin{equation}
p^{(c)} = \frac{\exp\!\left(\text{logit}^{(c)} - \max_j \text{logit}^{(j)}\right)}{\sum_{k=1}^{C} \exp\!\left(\text{logit}^{(k)} - \max_j \text{logit}^{(j)}\right)}
\end{equation}
where the subtraction of $\max_j \text{logit}^{(j)}$ prevents floating-point overflow in the exponential. The predicted class corresponds to $\arg\max_c\, p^{(c)}$, returned alongside the maximum probability as the classification confidence.

\subsection{Multi-Leaf Aggregation Strategies}
\label{subsec:aggregation_strategies}

Since vineyard images typically contain multiple leaves (5--15 visible leaves per $640 \times 480$ frame in field conditions), frame-level disease assessment requires aggregating predictions across all detected instances. The aggregation module implements three strategies, drawing from multi-instance learning (MIL) principles~\cite{dietterich1997solving}:

\textbf{1. Detection-Confidence Weighting (Baseline):} Each leaf's class probabilities are weighted by its detection confidence $c_{\text{det}}$:
\begin{equation}
p_{\text{frame}}^{(c)} = \frac{\sum_{i=1}^{N} c_{\text{det}}^{(i)} \cdot p_i^{(c)}}{\sum_{i=1}^{N} c_{\text{det}}^{(i)}}
\end{equation}
This baseline assumes higher-confidence detections correspond to better-localized leaves with more reliable classification features.

\textbf{2. Entropy-Based Uncertainty Weighting:} Shannon entropy~\cite{shannon1948mathematical} quantifies per-leaf classification uncertainty:
\begin{equation}
H_i = -\sum_{c=1}^{C} p_i^{(c)} \log p_i^{(c)}
\end{equation}
Leaves with low entropy (confident predictions) receive higher aggregation weight than ambiguous predictions. The normalized weight becomes:
\begin{equation}
w_i = c_{\text{det}}^{(i)} \cdot \left(1 - \frac{H_i}{\log C}\right)
\end{equation}
where $\log C$ normalizes entropy to $[0, 1]$. This approach downweights leaves near class boundaries (e.g., Black Rot--Esca confusion) exhibiting high classification uncertainty.

\textbf{3. Hybrid Spatial-Quality Weighting (Production):} Bounding box geometry provides additional quality indicators, combined with detection confidence and entropy:

\begin{itemize}
    \item \textbf{Relative size:} $r_{\text{size}} = A_{\text{bbox}} / A_{\text{frame}}$, penalizing excessively small (noise) or large (misdetections) boxes outside the range $[0.005, 0.5]$ of frame area.
    \item \textbf{Aspect ratio:} Scoring function centered at ideal ratio 1.0 with tolerance 0.5, penalizing elongated boxes corresponding to stems or trellis wires.
    \item \textbf{Centrality:} $d_{\text{center}} = \|(x_c, y_c) - (W/2, H/2)\|$, downweighting edge-positioned leaves that may be partially occluded or distorted by lens aberration.
\end{itemize}

The combined quality score $q_i \in [0, 1]$ multiplies the detection-entropy weight:
\begin{equation}
w_i^{\text{hybrid}} = c_{\text{det}}^{(i)} \cdot \left(1 - \frac{H_i}{\log C}\right) \cdot q_i
\end{equation}

The frame-level diagnosis selects the class $c^* = \arg\max_c\, p_{\text{frame}}^{(c)}$ with the aggregated confidence as the final output.


\section{Implementation Constraints and Design Trade-offs}
\label{sec:implementation_constraints}

\subsection{Build System and Model Embedding Constraints}
\label{subsec:build_constraints}

Embedding neural network models as RODATA in the firmware binary introduces build system constraints that affect the development workflow. Both \texttt{.espdl} model files are referenced via GCC inline assembly directives (\texttt{asm("\_binary\_custom\_detect\_espdl\_start")}) that resolve at link time. Any modification to model files --- including re-quantization with different calibration data, architecture changes, or weight updates --- requires a full clean build (\texttt{idf.py fullclean}) to ensure the linker symbol references point to the updated model binaries. Incremental builds may silently retain stale model data, leading to initialization failures or incorrect inference results that are difficult to diagnose.

The production v29 firmware binary occupies 4,356,032 bytes (4.15~MB), comprising 2,971,836 bytes of memory-mapped model data (segment 0, mapped to virtual address \texttt{0x3c140020}) and 1,268,968 bytes of application code (segment 3). This binary size fits within the 5~MB factory partition but leaves limited margin for firmware expansion, such as the addition of communication stacks (WiFi, LoRaWAN) or additional model variants.

Binary size verification serves as a critical deployment check: a firmware binary substantially smaller than expected (e.g., 1.28~MB instead of 4.15~MB) indicates that model embedding failed silently, and the system will crash at boot with null pointer access when attempting to load models from \texttt{.rodata}.

\subsection{Class Index Mapping and Framework Compatibility}
\label{subsec:class_mapping}

Deploying PyTorch-trained models on the ESP-DL runtime introduces a cross-framework interoperability challenge in class label ordering. PyTorch's \texttt{ImageFolder} dataset loader assigns class indices based on alphabetical directory sorting: \texttt{Black\_rot}$\rightarrow$0, \texttt{Esca}$\rightarrow$1, \texttt{Healthy}$\rightarrow$2, \texttt{Leaf\_blight}$\rightarrow$3. The class name array in the ESP32 firmware must replicate this exact ordering for correct label assignment from model output indices.

A related challenge concerns the dequantization interface asymmetry in the ESP-DL framework. The detection post-processor (\texttt{ESPDetPostProcessor}) handles INT8-to-float conversion internally, transparently returning floating-point confidence scores and bounding box coordinates. The classification model, however, outputs raw INT8 logit values through the \texttt{dl::Model} interface, requiring explicit dequantization via multiplication by $2^{\text{exponent}}$ (Equation~\ref{eq:dequant}) before softmax probability computation. Without this step, raw INT8 values in the range $[-128, 127]$ produce extreme softmax outputs (near-100\% for a single class) due to the exponential amplification of large magnitude differences.

As a concrete illustration, consider an output tensor with exponent $e = -7$ (scale $= 0.0078125$). Raw INT8 logits $[23, -15, 87, 12]$ should dequantize to $[0.180, -0.117, 0.680, 0.094]$, yielding a meaningful probability distribution via softmax. Without dequantization, softmax of $[23, -15, 87, 12]$ produces $[\approx 0\%, \approx 0\%, \approx 100\%, \approx 0\%]$, masking any meaningful class discrimination. These constraints underscore the importance of end-to-end integration testing when deploying quantized models across heterogeneous frameworks.

Having described the complete deployment pipeline and the practical challenges encountered, the next chapter evaluates the system's performance in terms of inference latency, classification accuracy, memory utilization, power consumption, and comparison with alternative edge computing platforms.
